\documentclass[12pt,a4paper]{article}

\input{/home/tabaka/Documents/GitHub/Defs/defs.tex}

%\pagestyle{empty} 			%usuwa nr strony

\setcounter{zadaniaRozdzial}{3}


\begin{document}
	
	\begin{center}
		\LARGE 2.3 Funkcja kwadratowa z parametrem
	\end{center}
	
	\Komentarz{Zadania z parametrem są bardzo schematyczne, można ich sposób rozwiązywania rozpisać w następujący sposób:
	\begin{enumerate}[a)]
		\item Rozpisujemy warunki
		\item $\Delta > 0$
		\item $a\neq0$
		\item $x_1\cdot x_2 \dots 0, \qquad x_1+x_2 \dots 0$
		\item Inny warunek zadania
	\end{enumerate}
	\begin{enumerate}[\text{Ad} a)]
		\item A następnie je po koleji rozwiązujemy
		\item Warunek by instaniały przynajmniej dwa miejsca zerowe
		\item Warunek by to była funkcja kwadratowa \textcolor{red}{(o tym się bardzo często zapomina)}
		\item Wzory Viete'a - jeśli pytają się o jakieś warunki co do pierwiastków f. kwadratowej (np. oba są dodatnie, to $x_1\cdot x_2>0 \wedge x_1+x_2>0$)
		\item Najczęściej trzeba rozpisać ze wzorów Viete'a (Szczególnie jak zadanie jest za 5-6 pkt)
	\end{enumerate}
	}
	
	\ZbiorZadanie{Wyznacz wartości parametru $k\in\mathbb{R}$, dla którego
	\begin{enumerate}[a)]
		\item $8(k^2-1)x^2+(8k-8)x+1=0$ nie ma rozwiązań,
		\item $(11-4m)x^2+(2m-4)x-1=0$ ma dokładnie jedno rozwiązanie,
		\item $(k^2-3k+2)x^2+3(k-2)x+4\frac{1}{2}=0$ ma dwa rozwiązania.
	\end{enumerate}}

	\ZbiorZadanie{Wykazać, że dla każdej liczby rzeczywistej $x$ i każdej liczby rzeczywistej $m$ prawdziwa jest nierówność 
		$$4x^2-2mx + m^2 \geq 6x+3m-9.$$}
		
	\ZbiorZadanie{Wyznaczyć wszystkie wartości parametru $k\in \mathbb{R}$ dla których równanie
	$$x^2-2kx+4-k^2=0$$ ma dwa pierwiastki dodatnie.}
	
	\ZbiorZadanie{Wyznaczyć wszystkie wartości parametru $j\in \mathbb{R}$ dla których równanie
		$$x^2+5jx+4j^2-3j=0$$ ma dwa pierwiastki ujemne.}
	\newpage	
	\ZbiorZadanie{Wyznaczyć, dla jakich wartości parametru $k\in \mathbb{R}$, różne rozwiązania $s_1,x_2$ równania
	$$kx^2-(k+1)x-2k+3=0$$
	spełniają warunek $\frac{1}{x_1}+\frac{1}{x_2}=k+1$.}
	
	\ZbiorZadanie{Wykazać, że jeśli różne rozwiązania $x_1,x_2$ równania $$x^2-px+p+8=0$$ spełniają warunek $$x_1^2x_2+x_1x_2^2\geq x_1+x_2-6,$$ to $p\in (-\infty,-6\rangle\cup(8,\infty)$.}
	
	\ZbiorZadanie{Dane jest równanie $$(2m+1)x^2-(m+3)x+2m+1=0$$. Wyznaczyć te wartości parametru $m\in\mathbb{R}$, dla którego suma odwrotności różnych pierwiastków tego równania jest większa od 1.}
	
	\ZbiorZadanie{Wyznaczyć wartości parametru $a\in\mathbb{R}$, dla których różne pierwiastki $x_1,x_2$ równiania 
	$$x^2-3x-a+1=0$$
	spełniają warunek $3x_1-2x_2=4$.}
	
	\ZbiorZadanie{Wyznaczyć, dla jakich wartości parametru $k\in\mathbb{R}$ równanie
	$$x^2-2x-\frac{p-5}{p+3}$$
	ma dwa pierwiastki jednakowych znaków, których suma kwadratów jest nie mniejsza od $3$.}
	
	\ZbiorZadanie{Dane są dwie funkcje
	$$f(x)=2x^2+x-m \qquad g(x)=mx^2-2mx+3.$$
	Wyznaczyć dla jakich wartości parametru $m\in\mathbb{R}$ wykresy funkcji $f$ i $g$ przecinają sie w dwóch punktach, których odcięte mają różne znaki.}
	
	\ZbiorZadanie{Wyznaczyć dla jakich wartości parametru $k\in\mathbb{R}$ poniższa nierówność nie ma rozwiązań.
	$$(k-2)x^2+(k-2)x+k-1<0$$}
	
	\ZbiorZadanie{Dla jakich wartości parametru $m\in\mathbb{R}$ zbiór rozwiązań nierówności $$x^2-3x+2<0$$
	jest zawarty w zbiorze rozwiązań nierówności $$mx^2-(3m+1)x+3>0?$$}
	
	\newpage
	
	\begin{center}
		\LARGE Zadania uzupełniające
	\end{center}
	\vspace{1.5cm}
	
	\ZbiorZadanie{Wykazać, że dla każdej liczby rzeczywistej $x$ i każdej liczby rzeczywistej $m$ prawdziwa jest nierówność 
		$$2x^2-mx + 2 > 3x-m^2$$}

	\ZbiorZadanie{Wyznaczyć wszystkie wartości parametru $k\in \mathbb{R}$ dla których równanie
		$$x^2-(2k-4)x+k^2-2k-k=0$$ ma dwa pierwiastki dodatnie.}
	
	\ZbiorZadanie{Wyznaczyć wszystkie wartości parametru $k\in \mathbb{R}$ dla których równanie
		$$x^2+jx-j+3=0$$ ma dwa pierwiastki ujemne.}
	
	\ZbiorZadanie{Wyznaczyć wszystkie wartości parametru $k\in \mathbb{R}$ dla których równanie
		$$x^2+2(k+4)x+k^2-2k=0=0$$ 
		ma dwa pierwiastki różnego znaku. (Rozważyć dwa przypadki, kiedy jest dodatni i ujemny pierwiastek oraz jeden z pierwiastków jest równy 0, które jest neutralne)}
		
	\ZbiorZadanie{Wyznaczyć wartość parametru $k\in\mathbb{R}$, dla których różne rozwiązania $x_1$ $x_2$ funkcji kwadratowej 
	$$x^2-(2k+1)x+k^2+2=0$$ spełniają warunek $x_1=2x_2$.}
	
	\ZbiorZadanie{Dnaje jest równanie $$x^2-mx+2=0.$$ Wyznaczyć, dla jakich wartości parametru $m\in\mathbb{R}$ wszystkie rozwiązania tego równania należą do przedziału $\langle0,3\rangle$.}
	
	
\end{document}